The first realization of the International Lunar Reference Frame (ILRF) is
limited at the decimetre level by a strong correlation ($\corr \approx -0.97$)
between the lunocentre $X$-translation and the frame scale, an origin error of
$\approx\!15$~cm traced to the near-side, near-equatorial geometry of the lunar
laser-ranging (LLR) retroreflector network. We ask a design question behind that
empirical result: \emph{which} additional observations resolve the degeneracy,
and at what cost. We answer it in three parts. First, we prove a null-space
classification theorem for the internal-ranging datum problem: a single range
observation contributes rank one to the datum Fisher information and leaves a
six-dimensional null space; the origin--scale pair is near-null for a near-side
cluster under a fixed body orientation; and physical libration lifts the datum
defect to zero while leaving the pair near-degenerate. Second, we build a unified
multi-technique datum-information substrate over the full seven-parameter Helmert
datum---Earth--reflector LLR, two-station lunar VLBI differential delay, and
orbiter-to-beacon range---and use it to show that the lunocentre-$X$/scale
degeneracy is a \emph{radial} ambiguity: radial (depth-diverse) ranging, such as
far-side orbiter tracking, breaks it directly, whereas transverse VLBI helps only
indirectly because its information lands in the transverse translation. Third, we
cast frame realization as cost-aware, gauge-constrained optimal experiment design
and compute a cost/degeneracy Pareto frontier; a far-side orbiter range delivers
$\approx\!8\times$ the degeneracy-breaking gain of a limb VLBI baseline for
comparable exposure. The decomposition linear algebra is validated to machine
precision against an independent SciPy/NumPy reference; all reported lunar
magnitudes are explicitly \emph{modelled} (representative geometry, noise, and
cost), so the deliverable is the \emph{structure and ordering} of the trade, not a
mission recommendation. The classification and the radial-diversity design law
give ILRF realization and lunar-PNT programmes (ESA Moonlight/LCNS, NASA LunaNet)
a principled basis for prioritizing tracking geometry over tracking volume. The
substrate is released open source and every reported number is reproducible from a
single command.
